\section{A primer for readers new to the algebra}\label{ch:primer}
\sectionbrief
{only linear algebra, polynomials and linear ODEs with constant coefficients.}
{the whole theory speaks the language of rings and ideals. Every notion used later is introduced here from a familiar starting point: integers first, then ODEs in one variable, then two variables.}
{nothing later depends on reading this chapter, but every later chapter uses its vocabulary. It covers ideals and annihilators (\S\ref{sec:pr-ideal}), exponential-polynomials (\S\ref{sec:pr-ep}), the Chinese remainder theorem (\S\ref{sec:pr-crt}) and the companion matrix (\S\ref{sec:pr-comp}).}

\subsection{Ideals and annihilators, starting from integers}\label{sec:pr-ideal}

\subsubsection{Integers}

The multiples of $6$ form the set
\[
(6)=\{\dots,-12,-6,0,6,12,\dots\}.
\]
This set has two properties. The sum of two multiples of $6$ is a multiple of $6$. A multiple of $6$ times any integer is again a multiple of $6$. A subset of a ring with these two properties is called an \emph{ideal}.

Now take the ideal generated by $4$ and $6$, that is, all combinations $4a+6b$ with integers $a,b$:
\[
(4,6)=\{4a+6b:\ a,b\in\Z\}.
\]
Since $2=6-4$ lies in it, and every element is even, we get $(4,6)=(2)$. The ideal ``remembers'' only the greatest common divisor of its generators. In $\Z$ every ideal is generated by one number.

\subsubsection{Operators in one variable}

Replace integers by differential operators with constant coefficients in one variable $x$, and replace ``is a multiple of $6$'' by ``kills the function $u$''.

\begin{definition}[{Annihilator (one variable)}]
For a function $u(x)$, the annihilator of $u$ is the set of all polynomials $Q$ such that the operator $Q(d/dx)$ kills $u$:
\[
\Ann(u)=\Big\{Q\in\C[\xi]:\ Q\Big(\frac{d}{dx}\Big)u=0\Big\}.
\]
\end{definition}

It is an ideal. If $Q$ kills $u$ and $S$ is any polynomial, then $SQ$ kills $u$ too, because $S(d/dx)$ applied to $0$ is $0$. Sums of annihilators are annihilators.

\begin{example}[{Annihilators of familiar functions}]
Write $\partial=d/dx$. The operator $\partial-1$ kills $e^x$, and so does every multiple of it. The basic cases are:
\begin{align*}
\Ann(e^x)&=(\partial-1),\\
\Ann(\sin x)&=(\partial^2+1),\\
\Ann(xe^x)&=\big((\partial-1)^2\big),\\
\Ann(e^x+e^{2x})&=\big((\partial-1)(\partial-2)\big).
\end{align*}
The function $e^{x^2}$ has annihilator $\{0\}$: no nonzero polynomial in $\partial$ kills it, because each derivative $\partial^ke^{x^2}$ equals a polynomial of degree exactly $k$ times $e^{x^2}$.
\end{example}

In one variable every ideal has a single generator, exactly as in $\Z$. The generator of $\Ann(u)$ is the \emph{minimal annihilator} of $u$: the ODE of lowest order that $u$ satisfies.

\subsubsection{Two variables: the new feature}

With two variables $t$ and $x$ we have two derivatives. An operator is a polynomial in both. Here an ideal may need \emph{several} generators.

\begin{example}[{An ideal with two generators}]
Let $u=e^{2t+3x}$. Then
\[
(\partial_t-2)u=0\qquad\text{and}\qquad(\partial_x-3)u=0 .
\]
Every operator of the form $A\,(\partial_t-2)+B\,(\partial_x-3)$ kills $u$, and these are all of them. So
\[
\Ann(u)=(\partial_t-2,\ \partial_x-3).
\]
No single polynomial generates this ideal: a polynomial in the ideal vanishes at the point $(2,3)$, and a single polynomial vanishes on a whole curve, not at a single point.
\end{example}

This is the source of every subtlety in the theory. A solution of a PDE has an annihilator that may contain a ``curve part'' (one polynomial) and a ``point part'' (finitely many points). Lemma~\ref{f:lem:gcd} of the next chapter makes this precise.

\subsection{Exponential-polynomials}\label{sec:pr-ep}

An \emph{exponential-polynomial} is a finite sum of polynomials times exponentials:
\[
w(t,x)=\sum_{j=1}^{r}p_j(t,x)\,e^{\tau_jt+\xi_jx}.
\]
These are exactly the solutions of ODEs with constant coefficients, in each variable. They are the ``trivial'' solutions of the theory: they live on finitely many points of the characteristic variety.

The key property is a finite-dimension test. Differentiate $w$ as often as you like: you always stay inside a finite-dimensional space.

\begin{example}[{Which functions are exponential-polynomials?}]
\begin{center}\small
\begin{tabular}{@{}L{4.4cm}L{1.6cm}L{8.0cm}@{}}
\toprule
Function & In $\EP$? & Span of all derivatives\\\midrule
$xe^{x}$ & yes & $\spn\{xe^x,\ e^x\}$, dimension $2$\\
$x^2e^{3t}+\cos x$ & yes & finite: $\cos x=\tfrac12(e^{ix}+e^{-ix})$\\
$e^{-x^2}$ & no & $\partial_x^ke^{-x^2}=(\text{degree }k)\cdot e^{-x^2}$: infinite dimension\\
$1/(1+x^2)$ & no & denominators $(1+x^2)^{k+1}$ grow: infinite dimension\\
$f(x+t)$, $f$ a Gaussian & no & same as the Gaussian\\
\bottomrule
\end{tabular}
\end{center}
The test ``is the span of the derivatives finite-dimensional?'' is Lemma~\ref{f:lem:ep}.
\end{example}

\subsection{The Chinese remainder theorem}\label{sec:pr-crt}

\subsubsection{Integers}

Solve the system
\[
x\equiv1\pmod 2,\qquad x\equiv2\pmod 3 .
\]
The unique solution modulo $6$ is $x\equiv5$. In general, if $a$ and $b$ are coprime, a remainder modulo $a$ and a remainder modulo $b$ determine a unique remainder modulo $ab$. In the language of rings:
\[
\Z/(ab)\;\cong\;\Z/(a)\times\Z/(b).
\]

\subsubsection{ODEs in one variable}

Let $p(\partial)=(\partial-1)(\partial-2)$. Its solutions are $u=ae^x+be^{2x}$. The two pieces can be extracted by \emph{projector operators}:
\begin{align*}
u_1&=\frac{\partial-2}{1-2}\,u=a\,e^{x},\\[3pt]
u_2&=\frac{\partial-1}{2-1}\,u=b\,e^{2x}.
\end{align*}
Check the first line: $(\partial-2)(ae^x+be^{2x})=-ae^x$, and dividing by $1-2=-1$ gives $ae^x$. The sum is direct:
\[
\ker p(\partial)=\ker(\partial-1)\oplus\ker(\partial-2).
\]
This works because $\partial-1$ and $\partial-2$ are coprime: there are polynomials $a,b$ with $a(\partial-1)+b(\partial-2)=1$ (here $a=1$, $b=-1$).

\subsubsection{Two variables: what changes}

For the wave operator $P=\tau^2-\xi^2=(\tau-\xi)(\tau+\xi)$ the two factors are coprime, but no combination of them equals $1$. Indeed
\[
(\tau-\xi,\ \tau+\xi)=(\tau,\ \xi),
\]
the ideal of the single point $(0,0)$ where the two lines meet. The quotient
\[
\ring/(\tau,\xi)\cong\C
\]
has dimension $1$. This number is the ambiguity of the splitting: the d'Alembert components $f(x+t)$ and $g(x-t)$ are defined only up to a constant moved from one to the other. Theorem~\ref{f:thm:factor} says that this finite overlap is the \emph{only} failure of the integer picture.

\subsection{The companion matrix}\label{sec:pr-comp}

A linear ODE of order $m$ becomes a first-order system for the vector of derivatives. For
\[
u''-3u'+2u=0
\]
put $X=(u,u')^{\top}$. Then $X'=CX$ with
\[
C=\begin{pmatrix}0&1\\-2&3\end{pmatrix}.
\]
The characteristic polynomial of $C$ is $\lambda^2-3\lambda+2$, the symbol of the ODE, with roots $1$ and $2$. The Lagrange projectors are
\[
\Pi_1=\frac{C-2\,\Id}{1-2},\qquad \Pi_2=\frac{C-\Id}{2-1},
\]
and the matrix exponential splits into modes:
\[
e^{xC}=e^{x}\,\Pi_1+e^{2x}\,\Pi_2 .
\]
This is the same formula as the projectors of the previous section, written for matrices. It reappears as the Lagrange formula~\eqref{f:eq:lagrange} for evolution factors.

\begin{keypoint}
Four ideas carry the theory. An \emph{ideal} records all equations a function satisfies. Its \emph{gcd} separates a curve part from finitely many points. \emph{Exponential-polynomials} are the functions with finitely many derivatives up to linear dependence. The \emph{Chinese remainder theorem} splits a solution into pieces, up to a finite overlap.
\end{keypoint}

